\section{Conceptos fundamentales}

\resumebox{definicion}{Fuerza magnética del protón}{\footnotesize
    \textbf{Cromatografía:} Técnica de separación que distribuye los componentes de una mezcla entre una fase estacionaria y una fase móvil.
}

\subsection{Tipos de cromatografía}

\begin{itemize}
    \item \textbf{Cromatografía de gases (CG):} La fase móvil es un gas. Se usa para compuestos volátiles.
    \item \textbf{Cromatografía de líquidos de alta resolución (HPLC):} La fase móvil es un líquido a alta presión. Se usa para una amplia gama de compuestos.
\end{itemize}

\section{Componentes de un cromatógrafo}

\begin{itemize}
    \item \textbf{Fase estacionaria:} Sustancia sólida o líquida inmóvil que retiene a los componentes de la muestra.
    \item \textbf{Fase móvil:} Líquido o gas que arrastra a los componentes.
    \item \textbf{Columna:} Tubo donde se encuentra la fase estacionaria.
    \item \textbf{Bomba (HPLC) o gas de arrastre (GC):} Impulsa la fase móvil.
    \item \textbf{Inyector:} Introduce la muestra en la columna.
    \item \textbf{Detector:} Mide la concentración de los compuestos al eluir de la columna.
\end{itemize}

\section{Detectores}
\begin{itemize}
    \item \textbf{Para CG:} Detector de ionización de flama (FID) y espectrometría de masas (MS).
    \item \textbf{Para HPLC:} UV-Vis, índice de refracción (RI), fluorescencia, y espectrometría de masas (MS).
\end{itemize}

\section{Conceptos clave}

\resumebox{ejercicio}{Fuerza magnética del protón}{\footnotesize
    \textbf{Tiempo de retención ($t_R$):} Tiempo que tarda un compuesto en eluir de la columna. Es característico de cada compuesto en condiciones específicas.
    \textbf{Tiempo de detección:} Tiempo en el que se registra la señal del compuesto en el detector.
}

\section{Cuantificación}

\begin{enumerate}
    \item \textbf{Curvas de calibración:} Con estándares de concentración conocida.
    \item \textbf{Estándar interno:} Se añade una cantidad conocida de un compuesto similar al analito para corregir variaciones en el proceso.
\end{enumerate}

\resumebox{ejercicio}{Fuerza magnética del protón}{\footnotesize
    \textbf{Interpretación de un cromatograma:} Un cromatograma de HPLC muestra un pico a los 3.2 minutos para una muestra de café y a los 3.2 minutos para un estándar de cafeína. ¿Qué puedes concluir? ¿Cómo cuantificarías la cafeína en el café?
}

\resumebox{dato}{Fuerza magnética del protón}{\footnotesize
    \textbf{Análisis de alimentos y bebidas:} La cromatografía es la técnica por excelencia en el control de calidad de alimentos. Se usa para determinar la presencia de aditivos, conservadores, pesticidas, micotoxinas y para verificar la autenticidad de un producto (ej. presencia de aceites de oliva adulterados).
}